\section{Related Work}


%Our method harnesses multiple LLM agents. 
Related works can be categorized into three parts: 1) LLM self-ensemble \cite{cotsc}, which attempts to harness multiple outputs from homogeneous LLMs to assemble the final answer; 2) heterogeneous LLM ensemble, which focuses on combining heterogeneous LLMs through supervised learning to improve performance across various downstream applications; and 3) multiple LLM agents collaboration, which improves performance through interactions among LLM agents. We discuss these works below. 



% The related work can be categorized into three parts: LLM self-ensemble, which attempts to harness multiple outputs from homogeneous LLMs to assemble the final answer \cite{reasoning_survey}; heterogeneous LLM ensemble, which focuses on combining heterogeneous LLMs through supervised learning to improve performance across various downstream applications; and multiple LLM agent collaboration, which enhances performance through interactions among LLM agents. These works implicitly incorporate elements of the "More Agents" mechanism; however, a comprehensive study of it has not been conducted until now. Our approach is pioneering in its thorough investigation of the "More Agents" mechanism. Our insights have the potential to orthogonally enhance these works.


% CoT-SC \cite{cotsc} harnesses diverse chain-of-thought \cite{cot} prompts to elicit a variety of reasoning processes from a single LLM, while \citet{complexity-cot} generates more complex reasoning processes, and thereafter the results are subjected to majority voting to produce the final answer. Essentially, it can be regarded as a fusion of our method with the CoT. In contrast, we also explore the combination with other methods.
% A series of works \cite{cot-sc, complexcot, diverse, verifier, lamda, agreement} leverage CoT to improve the performance of LLMs, which only focus on the reasoning capacity. Our method not only improves the reasoning performance but also the generation performance. Additionally, our method can integrate with multiple prompt engineering methods (including CoT) and multiple LLM-agents collaboration frameworks.
% A series of works \cite{cotsc, complexity-cot, diverse, verifier, lamda, self-agreement} leverage an ensemble of multiple CoTs to improve the performance of LLMs, which only focus on the reasoning capacity. Our method not only improves the reasoning performance but also the generation performance. Additionally, our method can integrate with multiple prompt engineering methods (including CoT) and multiple LLM-agents collaboration frameworks.




%CoT-SC \cite{cotsc} propose a sample-and-marginalize method based on LLMs via CoTs \cite{cot} to enhance the performance, claims that its performance increases with the number of sample paths.
\textbf{LLM Self-Ensemble.} CoT-SC \cite{cotsc} harnesses diverse chain-of-thought \cite{cot} prompts to elicit a variety of reasoning processes from a single LLM and select the final answer through majority voting. \citet{complexity-cot, diverse, verifier, lamda, self-agreement} can be considered as the extensions of CoT-SC. These methods mainly focus on reasoning tasks and exclusively investigate the compatibility with CoT. In contrast, our method not only validates effectiveness in reasoning tasks but also in generation tasks. Moreover, our method is compatible with a broader range of methods, such as prompt engineering (including CoT) and multiple LLM agents collaboration. Very recently, \citet{llm-blending} proposes a method named Blended that utilizes multiple LLMs for chat scenarios. In contrast, Blended focuses on utilizing the power of multiple LLMs, whereas our focus is on the scaling trend of adding more LLMs. Also, Blended is only for limited chat scenarios evaluated via human annotations. Furthermore, we explore orthogonality with other methods.
% Very recently, \citet{llm-blending} presents a method using an ensemble of smaller LLMs for performance akin to a larger model, and solely in chat applications. Our method, however, examines performance scalability with additional LLM agents in diverse tasks, including reasoning and generation.

% \sout{Overall, these studies are aimed at developing methods that leverage multiple LLMs to boost performance within specific scenarios. Our intention is to employ the simplest way with multiple LLMs to assess the universality and extent to which scaling LLMs can enhance performance. By analyzing the underlying principles, we aim to delineate the potential scope of application. Our goal is to devise a widely usable, lightweight, plug-in method.}
% Very recently, \citet{llm-blending} discovers that an ensemble of smaller LLMs can collaboratively match the performance of a single larger model. However, their study is primarily within a single chat scenario. In contrast, our method is more versatile and can augment other methods.

% Very recently, \citet{llm-blending} proposes a method to ensemble multiple LLMs based on Bayesian statistical principles to enhance the user retention and engagement in conversational scenario which are not the scope of this work.





% Very recently, \citet{llm-blending} proposes a method to ensemble multiple LLMs based on Bayesian statistical principles, specifically focusing on chat scenarios. 
%However, our method proves effective across a broader range of scenarios. 
% Other studies focus on post-processing techniques by using human annotations to train either an additional verifier \cite{verifier,diverse}, a re-ranker for response filtering \cite{lamda}, or prompting a language model to select the most consistent answer among the sampled reasoning paths \cite{self-agreement}. 
% By comparison, we report a finding that simply adding more instantiated LLM agents is what you need to obtain a better LLM performance in processing complex tasks, which is technically orthogonal to these methods. 
%These works can be viewed as post-processing applied to our method. 
%In practice, our method can orthogonally improve the aforementioned methods. 
 
\textbf{Heterogeneous LLM Ensemble. } \citet{fusellm} conducts a supervised LLM fusion framework to distill multiple heterogeneous LLMs into a single model and surpasses each of these LLMs. \citet{llm-blender} introduces a supervised ensembling framework based on multiple heterogeneous LLMs. \citet{frugalgpt} proposes a sequential inference method for LLMs that halts when the output quality is deemed adequate. \citet{foe} addresses the fusion-of-experts problem by integrating outputs from models with distinct knowledge domains through supervised learning. \citet{llms-routing} and \citet{rtte} select the most suitable LLM for new tasks by training a reward-guided router. These approaches primarily employ supervised learning, necessitating task-specific annotated data, and exhibit limited generalizability. In contrast, our method is unsupervised, without the need for additional training data.

% as detailed in Section \ref{sec:analysis}, we introduce  method that is both heterogeneous and complementary, requiring no additional training data and demonstrating commendable generalization capabilities.

\textbf{Multiple LLM Agents Collaboration. } \cite{debate,mad,debate3} explore various multiple LLM agents interaction architectures, with employing static debate-style engagements among LLMs for enhanced reasoning . \citet{dylan} enables agents to interact for multiple rounds in a dynamic architecture. \citet{camel,metagpt,autogen,agentverse,autoagent} offer several multi-agent frameworks that enable the development of LLM applications or enhance task-solving capabilities. However, these methods primarily focus on the interaction structures between LLM agents, rather than the relationship between the number of agents and performance. We also select representative methods \cite{debate,reflexion} to combine with our method, achieving further enhancements.


% We discern that these methodologies implicitly incorporate more agents mechanisms.
